\documentclass[ 12pt,a4paper,twocolumn,fleqn]{article}
\usepackage{graphicx}
\usepackage[a4paper,top=20mm, bottom=30mm, left=10mm, right=25mm]{geometry}
\usepackage{fancyhdr}
\usepackage{hyperref}
\usepackage{listings}
\usepackage[]{algorithm2e}
\usepackage{color}
\usepackage{fancybox}
\thisfancypage{%
  \setlength{\fboxsep}{20pt}\doublebox}{}
\pagestyle{fancy}
\usepackage{lineno}
\usepackage{xtab,booktabs}
\usepackage{setspace}
\usepackage{amsmath}
\usepackage{xcolor}
\usepackage{graphicx}
\usepackage{subcaption}
\usepackage{graphicx}
\usepackage{tikz}
\usepackage[compact]{titlesec}
\titlespacing{\section}{0pt}{*0}{*0}
\titlespacing{\subsection}{0pt}{*0}{*0}
\titlespacing{\subsubsection}{0pt}{*0}{*0}
\setlength\columnsep{25pt}
\makeatletter
\g@addto@macro{\normalsize}{%
\setlength{\abovedisplayskip}{0pt}
\setlength{\abovedisplayshortskip}{0pt}
\setlength{\belowdisplayskip}{0pt}
\setlength{\belowdisplayshortskip}{0pt}}
\makeatother
\mathindent=0.0pt
\usepackage{float}
\renewcommand{\baselinestretch}{1.5}
\begin{document}
\lhead{RL for Game Playing: AI agent training strategies}
\chead{}

\onecolumn
\begin{center}
\textcolor{brown}{\LARGE{DAYANANDA SAGAR UNIVERSITY}} \\
\includegraphics[scale=0.31]{media/SCHOOL OF ENGINEERING.png} \\

\large{\textbf{Bachelor of Technology}} \\
\large{in} \\
\large{Computer Science and Engineering} \\
\large{(ARTIFICIAL INTELLIGENCE AND MACHINE LEARNING)} \\

\includegraphics[width=0.18\textwidth]{media/Logo of AIML.jpeg.jpg}

\text{A Project Report On} \\
\smallskip
\textcolor{cyan}{\LARGE{REINFORCEMENT LEARNING FOR GAME PLAYING:
AI AGENT TRAINING STRATEGIES}} \\



\textit{Submitted By} \\
\textbf{Student 1 \space ENG00000000} \\
\textbf{Student 2 \space ENG00000000} \\
\textbf{Student 3 \space ENG00000000} \\
\textbf{Student 4 \space ENG00000000} \\
\textit{Under the guidance of} \\
\textbf{Prof/Dr. Professor name} \\
\text{Assistant Professor, CSE(AIML), DSU}\\
\large{\textbf{2023 - 2024}} \\
\textcolor{blue}{\large{Department of Computer Science and Engineering (AI \& ML)}} \\
\textcolor{blue}{\large{DAYANANDA SAGAR UNIVERSITY}} \\
\textcolor{blue}{\large{Bengaluru - 560068}} \\
\end{center}

\newpage
  \pagestyle{fancy}
\thisfancypage{%
  \setlength{\fboxsep}{20pt}\doublebox}{}
\begin{center}


\begin{figure}
  \centering
  \includegraphics[width=1\textwidth]{media/Untitled.png}
\end{figure}


\textcolor{brown}{\LARGE{Dayananda Sagar University}} \\

\footnotesize{Kudlu Gate, Hosur Road, Bengaluru - 560 068, Karnataka, India} \\
\begin{flushleft}
\textcolor{blue}{\LARGE{\textbf{Department of Computer Science \& Engineering}}} \\
\textcolor{blue}{\LARGE{\textbf{(Artificial Intelligence \& Machine Learning) 
}}} \\
\end{flushleft}
\Large{\underline{\textbf{CERTIFICATE}}} \\
  \end{center}
\normalsize
This is to certify that the project entitled \textbf{REINFORCEMENT LEARNING FOR GAME PLAYING:
AI AGENT TRAINING STRATEGIES} is a bonafide work carried out by \textbf{Student 1 (ENG00000000)}, \textbf{Student 2 (ENG00000000)}, \textbf{Student 3 (ENG00000000)} and \textbf{Student 4(ENG00000000)} in partial fulfillment for the award of degree in Bachelor of Technology in Computer Science and Engineering (Artificial Intelligence and Machine Learning), during
the year 2023-2024. \\
\\
\noindent
\begin{minipage}[t]{0.31\textwidth}
\small
  \textbf{Prof. Professor}\\
  Assistant Professor\\
  Dept. of CSE (AIML)\\
  School of Engineering\\
  Dayananda Sagar University\\
  \\
  Signature .........................\\
\end{minipage}%
\hfill
\begin{minipage}[t]{0.31\textwidth}
\small
  \textbf{Dr. Vinutha N}\\
  Project Co-ordinator\\
  Dept. of CSE (AIML)\\
  School of Engineering\\
  Dayananda Sagar University\\
  \\
  Signature .........................\\
\end{minipage}%
\hfill
\begin{minipage}[t]{0.34\textwidth}
\small
  \textbf{Dr. Jayavrinda Vrindavanam}\\
  Professor \& Chairperson\\
  Dept. of CSE (AIML)\\
  School of Engineering\\
  Dayananda Sagar University \\
  \\
  Signature .........................\\
\end{minipage}


\text{Name of the Examiners:}
\hfill
\text{Signature with date:} \\
\text{1 ...........................}
\hfill{ .............................} \\
\text{2 .............................} 
\hfill{ ............................} \\
\text{3 .............................} 
\hfill{ ............................} \\

% Latex code for adding double lined box for pages
% Should add this code for every page
\newpage
  \pagestyle{fancy}
\thisfancypage{%
  \setlength{\fboxsep}{20pt}\doublebox}{}
  
\begin{center}
\LARGE{{\underline{Acknowledgement}}} \\
\end{center}
\normalsize
It is a great pleasure for us to acknowledge the assistance and support of many individuals who have been responsible for the successful completion of this project work.
\\
First, we take this opportunity to express our sincere gratitude to \textbf{School of Engineering and Technology, Dayananda Sagar University} for providing us with a great opportunity to pursue our Bachelor’s degree in this institution.
\\
We would like to thank \textbf{Dr. Udaya Kumar Reddy K R}, Dean, School of Engineering and Technology, Dayananda Sagar University for his constant encouragement and expert advice.
\\
It is a matter of immense pleasure to express our sincere thanks to \textbf{Dr. Jayavrinda Vrindavanam}, Professor \& Department Chairperson, Computer Science and Engineering (Artificial Intelligence and Machine Learning), Dayananda Sagar University, for providing right academic guidance that made our task possible.
\\
We would like to thank our guide \textbf{Prof. Professor}, Assistant Professor, Dept. of Computer Science and Engineering, for sparing his valuable time to extend help in every step of our project work, which paved the way for smooth progress and fruitful culmination of the project.
\\
We would like to thank our Project Coordinator \textbf{Dr. Vinutha N} as well as all the staff members of Computer Science and Engineering (AIML) for their support.
\\
We are also grateful to our family and friends who provided us with every requirement throughout the course.
\\
We would like to thank one and all who directly or indirectly helped us in the Project work.


\begin{flushright}
\textbf{Student 1 \space ENG00000000} \\
\textbf{Student 2 \space ENG00000000} \\
\textbf{Student 3 \space ENG00000000} \\
\textbf{Student 4 \space ENG00000000} \\
\end{flushright}


\newpage
  \pagestyle{fancy}
\thisfancypage{%
  \setlength{\fboxsep}{20pt}\doublebox}{}
  
\setstretch{1.5}
\twocolumn[
\begin{@twocolumnfalse}
\title {REINFORCEMENT LEARNING FOR GAME PLAYING:
AI AGENT TRAINING STRATEGIES}
\date{}
\author{Student 1, Student 2, Student 3, Student 4}
\maketitle



{\bfseries\LARGE Abstract} \\
    
The project titled "Reinforcement Learning for Game Playing: AI Agent Training Strategies" is a comprehensive exploration of applying reinforcement learning techniques to train artificial intelligence (AI) agents for game-playing scenarios.
Leveraging the Unity ML-Agents framework, the project seeks to develop intelligent agents with the capacity to learn optimal strategies for effective navigation within dynamic game
environments, emphasizing the attainment of predefined goals and the avoidance of obstacles.
The study delves into a thorough investigation of diverse training strategies and algorithms to enhance the efficiency and adaptability of these AI agents.
Implementing and assessing various reinforcement learning
algorithms, including but not limited to Q-learning, Deep Q Networks (DQN), and Proximal
Policy Optimization (PPO). \\
Creating intricate game environments within Unity ML-Agents to simulate
realistic challenges for AI agents, fostering a dynamic learning environment. \\
Exploring a spectrum of training strategies, encompassing hyperparameter
tuning and optimization techniques, to maximize the learning efficiency of AI agents.
Goal Achievement and Obstacle Avoidance: Developing a framework to enable AI agents to
strategically achieve predefined goals while efficiently avoiding obstacles. \\
The outcomes of the project present promising capabilities of trained agents in achieving
predefined goals and navigating through obstacles.
The comparative analysis of reinforcement learning algorithms provides valuable insights into their respective strengths and limitations in the context of game playing.


\end{@twocolumnfalse}]
\modulolinenumbers[5]
\onecolumn
\newpage
  \pagestyle{fancy}
\thisfancypage{%
  \setlength{\fboxsep}{20pt}\doublebox}{}
\tableofcontents
\newpage
  \pagestyle{fancy}
\thisfancypage{%
  \setlength{\fboxsep}{20pt}\doublebox}{}
\newpage
  \pagestyle{fancy}
\thisfancypage{%
  \setlength{\fboxsep}{20pt}\doublebox}{}
  
\section{Introduction}
Reinforcement Learning (RL) \cite{b1} stands as a pivotal methodology in the realm of artificial
intelligence, particularly in the context of game playing. This paper delves into the myriad
strategies employed for training AI agents using reinforcement learning techniques within the
domain of game playing. The exploration spans a spectrum of crucial aspects, including diverse reinforcement learning algorithms, innovative training methodologies, and nuanced neural network architectures. The primary objective is to bolster the decision-making capabilities of AI agents, fostering adaptability to intricate gaming environments and ultimately achieving heightened performance levels.
Through a meticulous examination of recent advancements and insightful case studies, this paper not only provides a comprehensive overview of the current state of reinforcement learning for game playing but also identifies trends and challenges prevalent in this rapidly evolving field.  Furthermore, the paper analyzes the interplay between exploration and exploitation \cite{b2}, emphasizing the delicate balance required for effective learning. It also explores transfer learning techniques, where knowledge gained in one gaming context can be leveraged for accelerated learning in a different environment.
This introduction explores the landscape of RL for game playing, emphasizing the evolving strategies employed to train AI agents effectively and efficiently, contributing to the broader field of artificial intelligence and computational gaming.

\subsection{Scope}

\clearpage

\newpage
 \pagestyle{fancy}
\thisfancypage{%
  \setlength{\fboxsep}{20pt}\doublebox}{}

\section{Problem Definition}

\clearpage

\newpage
 \pagestyle{fancy}
\thisfancypage{%
  \setlength{\fboxsep}{20pt}\doublebox}{}

\section{Literature Survey}

\clearpage

\newpage
 \pagestyle{fancy}
\thisfancypage{%
  \setlength{\fboxsep}{20pt}\doublebox}{}

\section{Methodology}

\subsection{Data Collection}
Data collection is the systematic process of gathering and accumulating relevant information from various sources. This phase involves defining the scope of the data, selecting appropriate sources, and employing methods such as surveys, experiments, or utilizing existing databases. The quality and reliability of collected data significantly impact the success of subsequent analysis and modeling.
\\

\subsection{Data Pre-processing}
Data preprocessing is a crucial step in preparing raw data for analysis. It involves cleaning and transforming the data to enhance its quality and usability. Tasks include handling missing values, addressing outliers, scaling features, encoding categorical variables, and performing normalization. Effective data preprocessing ensures that the data is in a suitable format for modeling, improving the accuracy and interpretability of machine learning algorithms.
\\

\subsection{Model Implementation}
In our major project, a total of 10 classification models were implemented for better understanding of the dataset and the domain
\\
\subsubsection{Logistic Regression}

\clearpage

\newpage
 \pagestyle{fancy}
\thisfancypage{%
  \setlength{\fboxsep}{20pt}\doublebox}{}

\section{Requirements}

\subsection{Functional Requirements}
Requirement 1 \\
Requirement 2 \\
Requirement 3\\

\subsection{Non- Functional Requirements}
Requirement 1\\
Requirement 2\\
Requirement 3\\

\clearpage

\newpage
 \pagestyle{fancy}
\thisfancypage{%
  \setlength{\fboxsep}{20pt}\doublebox}{}

\section{Results \& Analysis}

\begin{figure}[h]
    \centering
    \includegraphics[width=0.5\textwidth]{media/sample.png}
    \caption{Your Image Caption}
    \label{fig:your_image}
\end{figure}

\clearpage

\newpage
 \pagestyle{fancy}
\thisfancypage{%
  \setlength{\fboxsep}{20pt}\doublebox}{}
  
\section{Conclusion \& Future work}

\clearpage

\newpage
 \pagestyle{fancy}
\thisfancypage{%
  \setlength{\fboxsep}{20pt}\doublebox}{}

\section{References}
% Include your references here
\begin{thebibliography}{9}

\bibitem{b1}
Chan, Stephanie CY, et al. "Measuring the reliability of reinforcement learning algorithms." arXiv preprint arXiv:1912.05663 (2019).

\bibitem{b2}
Louis, Ruwaid, and David Yu. "A study of the exploration/exploitation trade-off in reinforcement learning: Applied to autonomous driving." (2019). \\
\url{https://www.diva-portal.org/smash/get/diva2:1336430/FULLTEXT01.pdf}


\end{thebibliography}

\end{document}
